\documentclass[aps,prc,onecolumn,superscriptaddress,nofootinbib,10pt]{revtex4-2}
\usepackage{amsmath,amssymb}
\usepackage{booktabs}
\usepackage[colorlinks=true,linkcolor=blue,citecolor=blue,urlcolor=blue]{hyperref}

\newcommand{\nB}{n_B}
\newcommand{\nC}{n_C}
\newcommand{\nS}{n_S}
\newcommand{\muB}{\mu_B}
\newcommand{\muC}{\mu_C}
\newcommand{\muS}{\mu_S}
\newcommand{\munue}{\mu_{\nu_e}}
\newcommand{\mueff}[1]{\mu^{\mathrm{eff}}_{#1}}
\newcommand{\hc}{(\hbar c)^3}

\begin{document}

\title{vMIT: the vector-interaction MIT bag model for quark matter \\
       --- as implemented in \texttt{eos/vmit}}

\author{M.~Guerrini}
\affiliation{University of Ferrara}

\begin{abstract}
The \texttt{eos/vmit} subpackage implements the MIT bag model of
Chodos et al.~\cite{Chodos1974} augmented by a repulsive vector interaction
between quarks, in the form used by Gomes et al.~\cite{Gomes2019} and by
Constantinou et al.~\cite{Constantinou2021,Constantinou2023}. The model has
four parameters --- the bag constant $B$, the vector coupling
$a = g_V^2/m_V^2$, and the current quark masses --- and no scalar
self-consistency: the quark masses are fixed, so the only mean field to solve
for is the vector one, and its equation of motion is algebraic. This note
states the Lagrangian and the mean field, every parameter with its value, the
single-flavour thermodynamics at $T = 0$ and $T > 0$ in closed form, the
conserved-charge bookkeeping, the residual of every mode the package closes
row by row in the order the solver assembles it, every quantity a solved point
returns, and the numerics that decide whether a state has been reached. It is
written from the code, as its specification: every equation here is asserted by
\texttt{eos/vmit/verify/run\_full\_check.py} or by \texttt{test/vmit/}.
\end{abstract}

\maketitle

\section{Lagrangian and mean field}
\label{sec:lagrangian}

Three quark flavours $q \in \{u, d, s\}$, confined by a bag of constant energy
density $B$, interact through an isoscalar--vector field $V^\mu$:
%
\begin{equation}
\mathcal{L} = \sum_q \bar\psi_q \big[ \gamma_\mu \big( i\partial^\mu
      - g_V V^\mu \big) - m_q \big] \psi_q
      - \tfrac14 V_{\mu\nu} V^{\mu\nu} + \tfrac12 m_V^2\, V_\mu V^\mu
      - B ,
\label{eq:lagrangian}
\end{equation}
%
with $V_{\mu\nu} = \partial_\mu V_\nu - \partial_\nu V_\mu$. The quark masses
are the \emph{current} masses; unlike a Nambu--Jona-Lasinio model there is no
scalar condensate and no gap equation, so $m_q$ is a parameter and not a
solved quantity. One coupling $g_V$ serves all three flavours: the vector
interaction is flavour blind.

In uniform matter at rest only $V^0$ survives, and its equation of motion is
algebraic:
%
\begin{equation}
m_V^2\, V^0 = g_V \sum_q n_q
\qquad\Longrightarrow\qquad
V \;\equiv\; g_V V^0 \;=\; \frac{g_V^2}{m_V^2} \sum_q n_q
        \;=\; a\, (\hbar c) \sum_q n_q ,
\label{eq:vectorfield}
\end{equation}
%
which is the only self-consistency in the model. The coupling is carried as
the single combination $a = g_V^2/m_V^2$, quoted in $\mathrm{fm}^2$; with the
densities in $\mathrm{fm}^{-3}$ the product $a \sum_q n_q$ is an inverse
length, and the factor $\hbar c = 197.3269804\ \mathrm{MeV\,fm}$ converts it
to MeV. The two couplings are not separately identifiable at mean-field level
--- only the ratio $g_V^2/m_V^2$ enters Eq.~(\ref{eq:vectorfield}) and
everything downstream of it --- so the code never carries them apart, and an
inference run over the quark sector varies $a$, not $g_V$ and $m_V$.

Because $V$ couples to the quark \emph{number} density and is flavour blind,
it shifts all three chemical potentials by the same amount. The kinetic
(effective) potential that enters the Fermi integrals is
%
\begin{equation}
\mueff{q} = \mu_q - V ,
\label{eq:mueff}
\end{equation}
%
and the mean-field problem is the fixed point $n_q = n_q(\mueff{q}, T, m_q)$
with $V$ given by Eq.~(\ref{eq:vectorfield}). Note what
Eq.~(\ref{eq:mueff}) does \emph{not} contain: there is no rearrangement
self-energy, because $a$ is a constant and not a function of the density. A
density-dependent coupling would add one; this model has none, and
$\Sigma^R \equiv 0$ throughout.

Since the same $V$ is subtracted from all three potentials, the differences
are untouched:
%
\begin{equation}
\mueff{u} - \mueff{d} = \mu_u - \mu_d = \muC , \qquad
\mueff{s} - \mueff{d} = \mu_s - \mu_d = \muS ,
\end{equation}
%
so the isospin and strangeness structure of a state lives entirely in the
physical potentials and the vector field carries none of it. Only the baryon
direction is shifted, $\muB^{\rm eff} = \muB - 3V$.

\section{Parameters}
\label{sec:parameters}

Four numbers set the equation of state. There is no scalar sector and no
density dependence, which is what makes the model cheap enough to scan and why
$B$ and $a$ are the two axes a hybrid parameter study moves along. They are
held in a frozen dataclass \texttt{Parameters} and are always an
\emph{argument} to every entry point, never module state.

\begin{table}[h]
\centering
\begin{tabular}{llll}
\toprule
symbol & code name & value & meaning \\
\midrule
---       & \texttt{name} & \texttt{"vMIT\_default"} & label carried into
                                     table headers and figure legends \\
$m_u$     & \texttt{m\_u} & $5.0$~MeV   & current up-quark mass \\
$m_d$     & \texttt{m\_d} & $7.0$~MeV   & current down-quark mass \\
$m_s$     & \texttt{m\_s} & $150.0$~MeV & current strange-quark mass \\
$a$       & \texttt{a}    & $0.2$~$\mathrm{fm}^2$ & vector coupling
                                                    $g_V^2/m_V^2$ \\
$B^{1/4}$ & \texttt{B4}   & $180.0$~MeV & bag constant, as its fourth root \\
\bottomrule
\end{tabular}
\caption{The parameter set returned by \texttt{Parameters.default()}.}
\label{tab:parameters}
\end{table}

$B$ is stored as $B^{1/4}$ because that is the form it is quoted in; the
property \texttt{Parameters.B} returns $B = (B^{1/4})^4$ in $\mathrm{MeV}^4$,
and the single place that divides by $\hc$ to reach $\mathrm{MeV\,fm^{-3}}$ is
the bag term of Eq.~(\ref{eq:vectorbag}). At $B^{1/4} = 180$~MeV,
$B/\hc = 136.63\ \mathrm{MeV\,fm^{-3}}$.

\paragraph{What these numbers are fitted to, and what they are not.} The
masses are the current masses and are of the order of the Particle Data Group
values~\cite{ParticleDataGroup2024}; because there is no chiral condensate
here to dress them, they enter the Fermi integrals directly and $m_s$ is the
only one large enough to matter --- at $m_u = 5$~MeV and $\mueff{u} \gtrsim
300$~MeV the light flavours are ultra-relativistic to better than
$5\times10^{-4}$ in $n$. The pair $(B^{1/4}, a)$ is \emph{not} a published
fit. vMIT's parameters are what a hybrid study scans, and which pair is right
depends on the hadronic model the quark phase is paired with: $B$ sets where
deconfinement becomes energetically possible and $a$ sets how stiff the quark
branch is above it, which is how a hybrid star reaches two solar masses. The
values of Table~\ref{tab:parameters} are the working set of this repository
and the set \texttt{test/baseline} is frozen at;
\texttt{Parameters(B4=..., a=..., m\_s=...)} returns a parametrization with
any subset changed and the rest left at Table~\ref{tab:parameters}. There is
no scan driver in the library: parameters are arguments, so a sweep over
$(B^{1/4}, a, m_s)$ is a loop in the caller that builds one parametrization
per sample and calls the mixed engine
(\texttt{eos.mixed.eos\_table(..., vmit\_params=...)}) once on each.
The literature values these are taken near are $B^{1/4} \simeq 180$~MeV
and $a$ in the range $0$--$0.3\ \mathrm{fm}^2$ used by
Gomes et al.~\cite{Gomes2019} and Constantinou
et al.~\cite{Constantinou2021,Constantinou2023}, who write the same vector
term as $G_V$ or $g_V^2/m_V^2$.

\paragraph{Sensitivity, for orientation.} $a = 0$ removes the vector term
entirely and leaves the original bag model of Chodos
et al.~\cite{Chodos1974}; $B^{1/4} = 0$ removes the bag. Both limits are
reachable and one of them is a \texttt{verify/} check: with $a = 0$ and
$B = 0$ the solved state must equal the Fermi integrals evaluated at the
physical potentials themselves, asserted to $10^{-10}$
(\S\ref{sec:verify}).

\paragraph{Three routes to a parameter set.} Model parameters are always
arguments, so all three have to exist. \emph{By name:}
\texttt{Parameters.default()} is the working set above, and
\texttt{Parameters.named('vMIT\_default')} takes it by name. vMIT ships exactly
one set, so the map has a single entry; it exists so that a caller sweeping
parameter sets need not know which models happen to have more than one.
\emph{A new set:} every field carries a default, so
\texttt{Parameters(B4=..., a=...)} names only what changes; the dataclass is
frozen, so \texttt{dataclasses.replace} is how a set already in hand is
modified, and there is no setter and no mutating helper. \emph{From
nuclear-matter parameters:} no route, and none is missing --- vMIT has no
nuclear sector, so there is no \texttt{nmp.py} and nothing to invert.

\section{Single-flavour thermodynamics}
\label{sec:single}

Each flavour contributes as a free Fermi gas of mass $m_q$ and degeneracy
%
\begin{equation}
g = 2_{\mathrm{spin}} \times 3_{\mathrm{colour}} = 6 ,
\end{equation}
%
evaluated at $\mueff{q}$ with its antiparticles included. The code takes $g$
from the shared particle table rather than writing $6$.

\subsection{Finite temperature}

With $f(x) = [1 + e^{x/T}]^{-1}$ and $E_k = \sqrt{k^2 + m_q^2}$,
%
\begin{align}
n_q &= \frac{g}{2\pi^2 \hc} \int_0^\infty \! dk\, k^2
       \left[ f(E_k - \mueff{q}) - f(E_k + \mueff{q}) \right] ,
\label{eq:nq} \\
P_q &= \frac{g}{6\pi^2 \hc} \int_0^\infty \! dk\, \frac{k^4}{E_k}
       \left[ f(E_k - \mueff{q}) + f(E_k + \mueff{q}) \right] ,
\label{eq:Pq} \\
\varepsilon_q &= \frac{g}{2\pi^2 \hc} \int_0^\infty \! dk\, k^2 E_k
       \left[ f(E_k - \mueff{q}) + f(E_k + \mueff{q}) \right] .
\label{eq:epsq}
\end{align}
%
The factor $\hc$ is what carries $\mathrm{MeV}^3$ to $\mathrm{fm}^{-3}$ and
$\mathrm{MeV}^4$ to $\mathrm{MeV\,fm^{-3}}$; every public boundary of this
package is fm-based, and natural units never cross one. Note that $n_q$ is the
\emph{net} density, particles minus antiparticles, which is what is conserved
and what appears in Eq.~(\ref{eq:vectorfield}) and in the charges of
\S\ref{sec:charges}.

The entropy density is not integrated separately. It is obtained from the
three above through
%
\begin{equation}
s_q = \frac{\varepsilon_q + P_q - \mueff{q}\, n_q}{T} ,
\label{eq:sq}
\end{equation}
%
which is the Euler relation of a single ideal gas, and is exact rather than a
convenience: it is equal term by term to the entropy integral
%
\begin{equation}
s_q = \frac{g}{2\pi^2 \hc} \int_0^\infty \! dk\, k^2
      \sum_{\pm} \left[ \frac{x_\pm}{T} f(x_\pm)
                        + \ln\!\big( 1 + e^{-x_\pm/T} \big) \right] ,
      \qquad x_\pm = E_k \mp \mueff{q} .
\label{eq:sqint}
\end{equation}
%
Eq.~(\ref{eq:sq}) is written with $\mueff{q}$, not $\mu_q$: it is the identity
of the \emph{kinetic} gas, which is the gas the integrals describe. The step
back to the physical potentials is taken once, for the totals, in
\S\ref{sec:euler}.

These integrals are \emph{not} implemented in this subpackage. They come from
\texttt{eos.general.fermi\_integrals}, which evaluates them through the
Johns--Ellis--Lattimer analytic approximation~\cite{JohnsEllisLattimer1996},
uniformly valid from the degenerate to the non-degenerate limit and exact at
$T = 0$. They are written out here anyway, because a description that leaves
its ideal-gas integrals to a citation is not one a reader can reproduce the
model from.

\subsection{Zero temperature}

At $T = 0$ the antiparticle terms vanish, the occupations become step
functions, and the flavour is filled to the Fermi momentum
%
\begin{equation}
k_{F,q} = \begin{cases}
  \sqrt{\big(\mueff{q}\big)^2 - m_q^2} , & \mueff{q} > m_q , \\[2pt]
  0 , & \text{otherwise} ,
\end{cases}
\qquad E_{F,q} = \sqrt{k_{F,q}^2 + m_q^2} = \mueff{q} .
\label{eq:kF}
\end{equation}
%
The three quantities close in elementary functions,
%
\begin{align}
n_q &= \frac{g\, k_{F,q}^3}{6\pi^2 \hc} ,
\label{eq:nqT0} \\
\varepsilon_q &= \frac{g}{16\pi^2 \hc}
   \left[ k_{F,q} \big( 2 k_{F,q}^2 + m_q^2 \big) E_{F,q}
        - m_q^4 \ln \frac{k_{F,q} + E_{F,q}}{m_q} \right] ,
\label{eq:epsqT0} \\
P_q &= \frac{g}{48\pi^2 \hc}
   \left[ k_{F,q} \big( 2 k_{F,q}^2 - 3 m_q^2 \big) E_{F,q}
        + 3 m_q^4 \ln \frac{k_{F,q} + E_{F,q}}{m_q} \right] ,
\label{eq:PqT0}
\end{align}
%
and $s_q = 0$. Eqs.~(\ref{eq:nqT0})--(\ref{eq:PqT0}) satisfy
$\varepsilon_q + P_q = \mueff{q} n_q$, which is Eq.~(\ref{eq:sq}) at
$s_q = 0$; the $T \to 0$ limit of the finite-$T$ forms is the same, and
\texttt{eos.general.verify} asserts the agreement.

The threshold in Eq.~(\ref{eq:kF}) is what makes the strange flavour appear at
a finite density rather than at zero: with $m_s = 150$~MeV a state with
$\mueff{s} < 150$~MeV has $n_s = 0$ exactly at $T = 0$, and the onset is the
density at which $\mueff{s}$ crosses $m_s$. At $T > 0$ there is no threshold
--- the exponential tail populates the flavour at any $\mueff{s}$ --- which is
why the cold-start heuristic of \S\ref{sec:guesses} treats the strange
fraction as a function of both $T$ and $\nB$.

\subsection{The scalar density, and why it is not used}

The shared integral routine also returns a scalar density,
$\rho_{s,q} = (\varepsilon_q - 3 P_q)/m_q$, the trace identity of the
energy--momentum tensor for one gas. This subpackage discards it, and the
reason is structural rather than an oversight: a scalar density is what a
scalar field couples to, and vMIT has no scalar sector. The masses in
Eq.~(\ref{eq:lagrangian}) are current masses, not Dirac effective masses;
nothing in the model varies with $\rho_s$, and there is no gap equation for it
to source.

This matters for reading the code and the tables, because \texttt{n\_s} in
this package is \emph{not} a scalar density: it is the strange-quark number
density, the $s$ of $u, d, s$. The repository-wide identity
$n_s = (\varepsilon - 3P)/m^*$ has no meaning here --- there is no $m^*$ ---
and the symbol $n_s$ is free to mean the flavour density, which is what
Eqs.~(\ref{eq:charges}) consume.

\section{The vector field, the bag, and the totals}
\label{sec:totals}

Beyond the kinetic sum, two terms:
%
\begin{equation}
P_V = \varepsilon_V = \tfrac12\, a\, (\hbar c) \Big( \sum_q n_q \Big)^{\!2} ,
\qquad
P_B = -\,\frac{B}{\hc} , \qquad
\varepsilon_B = +\,\frac{B}{\hc} .
\label{eq:vectorbag}
\end{equation}
%
so the strongly-interacting sector totals
%
\begin{equation}
P^{\rm matter} = \sum_q P_q + P_V + P_B , \qquad
\varepsilon^{\rm matter} = \sum_q \varepsilon_q + \varepsilon_V
                           + \varepsilon_B , \qquad
s^{\rm matter} = \sum_q s_q .
\label{eq:mattertotals}
\end{equation}
%
\textbf{The two terms that differ between $P$ and $\varepsilon$ differ in
opposite ways, and this is the whole structure of the model.} The vector field
enters both with the \emph{same} sign, $\varepsilon_V = +P_V$: it is a
repulsion, it raises the pressure at fixed density, and it stiffens the quark
branch. The bag enters with \emph{opposite} signs,
$\varepsilon_B = -P_B = B/\hc$: it costs energy to make a bag, and it holds
the pressure negative --- the matter unbound --- until the kinetic and vector
terms have paid for it, which is what deconfinement at a finite density means
in this model. Neither term contributes to $s$: $V$ is a mean field with no
entropy of its own and $B$ is a constant. The \texttt{verify/} suite asserts
$\varepsilon_B = -P_B$ and $\varepsilon_V = +P_V$ to $10^{-14}$.

\subsection{Leptons, photons and what a point reports}

Eq.~(\ref{eq:mattertotals}) is the matter sector alone. What a solved point
reports adds, from \texttt{eos.general.thermodynamics\_leptons}, whichever of
these the mode and the flags call for:
%
\begin{itemize}
\item \textbf{electrons} (with positrons) at $\mu_e$, a Fermi gas of mass
      $m_e = 0.511$~MeV and degeneracy $2$, evaluated through the same
      Eqs.~(\ref{eq:nq})--(\ref{eq:sq}). Present wherever the mode has a
      lepton condition: both $\beta$-equilibrium modes always, and the two
      fixed-fraction modes when \texttt{leptons=True};
\item \textbf{electron neutrinos} (with antineutrinos) at $\munue$, massless
      and of degeneracy $1$, present only in
      \texttt{beta\_eq\_neutrino\_trapped};
\item \textbf{photons}, when \texttt{photons=True}: massless bosons at
      $\mu = 0$ and degeneracy $2$, so
      %
      \begin{equation}
      P_\gamma = \frac{\pi^2}{45}\, \frac{T^4}{\hc} ,\qquad
      \varepsilon_\gamma = 3 P_\gamma ,\qquad
      s_\gamma = \frac{4\pi^2}{45}\, \frac{T^3}{\hc} ,\qquad
      n_\gamma = \frac{2\zeta(3)}{\pi^2}\, \frac{T^3}{\hc} ,
      \label{eq:photons}
      \end{equation}
      %
      vanishing identically at $T = 0$.
\end{itemize}
%
None of these carries $B$, $C$ or $S$ in the sense of
Eq.~(\ref{eq:charges}) --- $\nC$ is non-leptonic by definition --- so they
enter $P$, $\varepsilon$ and $s$, and the neutrality condition, and nothing
else. The totals a point returns are
%
\begin{equation}
P = P^{\rm matter} + P_e + P_{\nu_e} + P_\gamma , \qquad
\varepsilon = \varepsilon^{\rm matter} + \varepsilon_e
              + \varepsilon_{\nu_e} + \varepsilon_\gamma , \qquad
s = s^{\rm matter} + s_e + s_{\nu_e} + s_\gamma ,
\label{eq:grandtotals}
\end{equation}
%
with $P_{\nu_e}$, $\varepsilon_{\nu_e}$, $s_{\nu_e}$ present only in the
trapped mode and the photon terms only when the flag is on, and with the free
energy density
%
\begin{equation}
f = \varepsilon - T s = -P + \sum_i \mu_i n_i ,
\end{equation}
%
the second equality being one of the \texttt{verify/} checks
(\S\ref{sec:verify}).

\subsection{The Euler relation}
\label{sec:euler}

With Eqs.~(\ref{eq:mueff}), (\ref{eq:vectorfield}) and (\ref{eq:vectorbag})
the identity
%
\begin{equation}
\varepsilon + P = T s + \sum_i \mu_i n_i
\label{eq:euler}
\end{equation}
%
holds \emph{identically}, not merely numerically, and the demonstration is
worth writing because it is what the vector term has to satisfy to be a mean
field rather than an added constant. For the matter sector: the $\pm B$ terms
cancel between $\varepsilon$ and $P$; the kinetic sector satisfies
Eq.~(\ref{eq:sq}) summed, $\sum_q (\varepsilon_q + P_q) = Ts + \sum_q
\mueff{q} n_q$; and what is left over is $2 P_V$, which by
Eqs.~(\ref{eq:vectorfield}) and (\ref{eq:vectorbag}) is
%
\begin{equation}
2 P_V = a (\hbar c) \Big( \sum_q n_q \Big)^{\!2} = V \sum_q n_q
      = \sum_q \big( \mu_q - \mueff{q} \big) n_q ,
\end{equation}
%
exactly the shift from the effective potentials to the physical ones. The
lepton and photon sectors each satisfy Eq.~(\ref{eq:euler}) on their own, so
they add without spoiling it. The \texttt{verify/} suite asserts
Eq.~(\ref{eq:euler}) at $10^{-8}$ relative in every mode; because the
cancellation above is exact, what that check actually tests is the Fermi
integrals and the assembly, with nothing left to hide an error in.

\section{Conserved charges}
\label{sec:charges}

Bookkeeping follows the conventions of the repository: $C$ is the electric
charge of strongly-interacting matter only (leptons excluded), and strangeness
counts $S = +1$ per $s$ quark, the opposite of the PDG sign. The quark quantum
numbers are
%
\begin{center}
\begin{tabular}{lccc}
\toprule
 & $B_q$ & $C_q$ & $S_q$ \\
\midrule
$u$ & $1/3$ & $+2/3$ & $0$ \\
$d$ & $1/3$ & $-1/3$ & $0$ \\
$s$ & $1/3$ & $-1/3$ & $+1$ \\
\bottomrule
\end{tabular}
\end{center}
%
from which
%
\begin{equation}
\nB = \frac{n_u + n_d + n_s}{3} , \qquad
\nC = \frac{2 n_u - n_d - n_s}{3} , \qquad
\nS = n_s ,
\label{eq:charges}
\end{equation}
%
and the fractions are $Y_C = \nC/\nB$ and $Y_S = \nS/\nB$, always relative to
$\nB$ and never to a total particle number. Species potentials are projections
of the conserved ones, $\mu_q = B_q \muB + C_q \muC + S_q \muS$, which for the
three flavours is a square system and inverts to
%
\begin{equation}
\muB = \mu_u + 2\mu_d , \qquad
\muC = \mu_u - \mu_d , \qquad
\muS = \mu_s - \mu_d .
\label{eq:potentials}
\end{equation}
%
The plus sign on $\muS$ in $\mu_s = \muB/3 - \muC/3 + \muS$ is the $S = +1$
convention; under the PDG sign it would be a minus. The sign of $\muC$ is
fixed by $\muC = \mu_p - \mu_n$ in the hadronic sector, and here reproduces it
through $\mu_p - \mu_n = (2\mu_u + \mu_d) - (\mu_u + 2\mu_d) = \mu_u - \mu_d$;
this is why $\beta$ equilibrium reads $\muC + \mu_e = 0$ in both sectors.

Neither Eq.~(\ref{eq:charges}) nor Eq.~(\ref{eq:potentials}) is written out in
this subpackage: both come from \texttt{eos.general.basis}, which builds them
from the quantum-number table above, so a quark phase and a hadronic phase
cannot drift apart in a hybrid construction.

\section{Equilibrium modes and their residuals}
\label{sec:modes}

A mode fixes the independent variables and supplies the conditions that close
the system. Every mode carries the same first four rows --- the three
flavour self-consistencies and the density --- and differs only in what comes
after them.

\subsection{The unknown vectors}
\label{sec:unknowns}

\textbf{There are two layouts, not one, and the difference is where $\mu_e$
sits.} In the $\beta$-equilibrium modes the lepton potential is a
\emph{potential} and travels with the other potentials, in slot four, ahead of
the densities. In the fixed-fraction modes it is appended \emph{after} the
densities, because it is present only conditionally and the layout must stay
the same length whether or not it is there. Writing either layout as though it
were universal --- as earlier drafts of this document did in both directions
--- makes every residual below index the wrong slot.

\begin{table}[h]
\centering
\begin{tabular}{lll}
\toprule
mode & \texttt{leptons} & unknown vector $x$ \\
\midrule
\texttt{beta\_eq\_neutrinoless}
  & --- & $(\mu_u, \mu_d, \mu_s, \mu_e,\; n_u, n_d, n_s)$ \hfill (7) \\
\texttt{beta\_eq\_neutrino\_trapped}
  & --- & $(\mu_u, \mu_d, \mu_s, \mu_e, \munue,\; n_u, n_d, n_s)$
          \hfill (8) \\
\texttt{fixed\_YC}
  & \texttt{True}  & $(\mu_u, \mu_d, \mu_s,\; n_u, n_d, n_s,\; \mu_e)$
                     \hfill (7) \\
\texttt{fixed\_YC}
  & \texttt{False} & $(\mu_u, \mu_d, \mu_s,\; n_u, n_d, n_s)$ \hfill (6) \\
\texttt{fixed\_YC\_YS}
  & \texttt{True}  & $(\mu_u, \mu_d, \mu_s,\; n_u, n_d, n_s,\; \mu_e)$
                     \hfill (7) \\
\texttt{fixed\_YC\_YS}
  & \texttt{False} & $(\mu_u, \mu_d, \mu_s,\; n_u, n_d, n_s)$ \hfill (6) \\
\bottomrule
\end{tabular}
\caption{The unknown vector of each mode, in the order the residual indexes
         it. In the $\beta$-equilibrium modes leptons are always present ---
         they are what the equilibrium is about --- so \texttt{leptons} does
         not apply.}
\label{tab:unknowns}
\end{table}

Carrying the densities $n_q$ as unknowns, rather than substituting
Eq.~(\ref{eq:vectorfield}) and solving for the potentials alone, is a
deliberate choice: it makes the residual \emph{polynomial} in the mean field
instead of nesting the Fermi integrals inside it. The cost is three extra
unknowns; the return is a smooth system a hybrid method closes in a few
iterations.

\subsection{The rows}

Throughout, $n_q^{\rm calc} \equiv n_q(\mueff{q}, T, m_q)$ is the density the
effective potentials produce through Eqs.~(\ref{eq:nq}) or (\ref{eq:nqT0}),
while $n_q$ without a superscript is the corresponding component of $x$, which
is what sources $V$ through Eq.~(\ref{eq:vectorfield}). The solver's job is to
make the two agree. Likewise $\nB(n_u,n_d,n_s)$, $\nC$ and $\nS$ are
Eqs.~(\ref{eq:charges}) evaluated on the $x$-components, not on the computed
ones, and $n_e(\mu_e)$, $n_{\nu_e}(\munue)$ are the lepton gases of
\S\ref{sec:totals}.

\paragraph{\texttt{beta\_eq\_neutrinoless}, conditions $(\nB, T)$.}
Seven rows, in the order assembled:
%
\begin{align}
R_1 &= n_u^{\rm calc} - n_u , &
R_2 &= n_d^{\rm calc} - n_d , &
R_3 &= n_s^{\rm calc} - n_s , \notag \\
R_4 &= \nB(n_u,n_d,n_s) - \nB , &
R_5 &= \nC(n_u,n_d,n_s) - n_e(\mu_e) , \notag \\
R_6 &= \mu_u + \mu_e - \mu_d , &
R_7 &= \mu_d - \mu_s . &&
\label{eq:res:beta}
\end{align}
%
$R_6$ is $\muC + \mu_e = 0$ written in flavour potentials --- the weak process
$d \leftrightarrow u + e + \bar\nu_e$ with $\munue = 0$ --- and $R_7$ is
$\muS = 0$, the strangeness-changing process $s \leftrightarrow d$. $R_5$ is
electric neutrality of the \emph{total} system, which is a separate statement
from $\nC$ itself and is imposed here because this mode says so.

\paragraph{\texttt{beta\_eq\_neutrino\_trapped}, conditions
$(\nB, Y_{L_e}, T)$.} Eight rows:
%
\begin{align}
R_1 &= n_u^{\rm calc} - n_u , &
R_2 &= n_d^{\rm calc} - n_d , &
R_3 &= n_s^{\rm calc} - n_s , \notag \\
R_4 &= \nB(n_u,n_d,n_s) - \nB , &
R_5 &= \nC(n_u,n_d,n_s) - n_e(\mu_e) , \notag \\
R_6 &= \mu_d - \mu_s , &
R_7 &= \mu_u + \mu_e - \mu_d - \munue , \notag \\
R_8 &= \frac{n_e(\mu_e) + n_{\nu_e}(\munue)}{\nB} - Y_{L_e} . &&
\label{eq:res:trapped}
\end{align}
%
\textbf{$R_6$ and $R_7$ swap places relative to
Eq.~(\ref{eq:res:beta}):} here the strangeness equality comes before the
$\beta$ condition. The rows are the same physics in both modes and the
ordering is immaterial to the solution, but it is not immaterial to a reader
matching this document against the residual, or to anyone writing a Jacobian
by hand, so it is stated as the code has it. $R_7$ is $\muC + \mu_e -
\munue = 0$: the neutrino potential is retained instead of being set to zero,
and $R_8$ fixes the electron-family lepton number in its place. The muon
family is not tracked (\S\ref{sec:deferred}), so $Y_{L_\mu}$ is refused rather
than ignored.

\paragraph{\texttt{fixed\_YC}, conditions $(\nB, Y_C, T)$.} With
\texttt{leptons=True}, seven rows:
%
\begin{align}
R_1 &= n_u^{\rm calc} - n_u , &
R_2 &= n_d^{\rm calc} - n_d , &
R_3 &= n_s^{\rm calc} - n_s , \notag \\
R_4 &= \nB(n_u,n_d,n_s) - \nB , &
R_5 &= \nC(n_u,n_d,n_s) - Y_C\, \nB , \notag \\
R_6 &= \mu_d - \mu_s , &
R_7 &= n_e(\mu_e) - \nC(n_u,n_d,n_s) . &&
\label{eq:res:yc}
\end{align}
%
With \texttt{leptons=False} the vector loses $\mu_e$ and the residual loses
$R_7$: six rows, $R_1$ through $R_6$, unchanged. There is no
$\beta$-equilibrium row --- $Y_C$ has replaced it --- but strangeness stays
equilibrated at $\muS = 0$ through $R_6$, which is what makes this a
three-flavour state at a prescribed charge rather than a two-parameter family.

\paragraph{\texttt{fixed\_YC\_YS}, conditions $(\nB, Y_C, Y_S, T)$.} With
\texttt{leptons=True}, seven rows:
%
\begin{align}
R_1 &= n_u^{\rm calc} - n_u , &
R_2 &= n_d^{\rm calc} - n_d , &
R_3 &= n_s^{\rm calc} - n_s , \notag \\
R_4 &= \nB(n_u,n_d,n_s) - \nB , &
R_5 &= \nC(n_u,n_d,n_s) - Y_C\, \nB , \notag \\
R_6 &= \nS(n_u,n_d,n_s) - Y_S\, \nB , &
R_7 &= n_e(\mu_e) - \nC(n_u,n_d,n_s) . &&
\label{eq:res:ycys}
\end{align}
%
With \texttt{leptons=False}, six rows, $R_1$ through $R_6$. \textbf{The
$\muS = 0$ row is gone}, replaced by $R_6$: once the strangeness fraction is
imposed, $\muS$ is an \emph{output} --- whatever potential the demanded amount
of strangeness costs --- and asking for both would over-determine the system.
This is the only mode in which $\muS \neq 0$ is expected, and it is what makes
$Y_C = 0.5$, $Y_S = 0$ the symmetric-nuclear-matter slice a heavy-ion
comparison wants.

\subsection{The \texttt{leptons} flag}

The flag is orthogonal to the mode and is an explicit argument, never one of
the conditions. With \texttt{leptons=True} an electron gas is added at
whatever $\mu_e$ makes the total system neutral --- that is exactly what the
row $R_7 = n_e - \nC$ says --- and it contributes to $\varepsilon$, $P$ and
$s$. With \texttt{leptons=False} the result is charged quark matter with no
leptons at all, which is what a mixed-phase construction needs for each pure
phase before imposing \emph{global} neutrality. In the two
$\beta$-equilibrium modes the flag has no meaning and is ignored: leptons are
what the equilibrium is about.

Photons are an independent flag, contribute to $\varepsilon$, $P$ and $s$
only, and carry no conserved charge. Wherever a temperature is accepted at a
\emph{point}, entropy per baryon $s/\nB$ is accepted in its place, through an
outer one-dimensional solve for $T$; the table driver does not yet take an
entropy axis (\S\ref{sec:deferred}).

\section{Numerics}
\label{sec:numerics}

\subsection{Judging a solution}

Convergence is judged on the residual vector, never on the root finder's own
termination report --- these are different questions, and a solver that
reports success has told you that its iteration stopped. The residual is made
dimensionless first, each row divided by the scale of the quantity it
balances:
%
\begin{equation}
\text{error} = \max_j \frac{|R_j|}{\sigma_j} ,
\qquad
\sigma_j = \begin{cases}
  \nB , & R_j \ \text{balances a density or a charge} , \\
  \max\big(|\muB|, 1\ \mathrm{MeV}\big) ,
        & R_j \ \text{equates two potentials} , \\
  1 , & R_j \ \text{is already dimensionless} ,
\end{cases}
\label{eq:scales}
\end{equation}
%
with $\muB = \mu_u + 2\mu_d$ evaluated at the current iterate, and a floor of
$1$~MeV so that a pathological iterate passing through $\muB = 0$ cannot
divide by zero. A state is accepted when $\text{error} < 10^{-10}$, which is
one number for the whole repository, so that ``converged'' means the same
thing in every model. Explicitly, per mode: in
Eq.~(\ref{eq:res:beta}) the scales are $(\nB,\nB,\nB,\nB,\nB,\muB,\muB)$; in
Eq.~(\ref{eq:res:trapped}), $(\nB,\nB,\nB,\nB,\nB,\muB,\muB,1)$ --- the
lepton-fraction row is already a ratio; in Eq.~(\ref{eq:res:yc}),
$(\nB,\nB,\nB,\nB,\nB,\muB,\nB)$; and in Eq.~(\ref{eq:res:ycys}) every row is
a density, so all seven scales are $\nB$.

This matters because the rows carry mixed units --- densities of order
$10^{-1}\ \mathrm{fm}^{-3}$, fractions of order unity, and equalities between
potentials of order $10^{3}$~MeV --- so a gate on the raw norm is dominated by
whichever row happens to be largest and accepts states that satisfy the others
only loosely.

\subsection{The solve}

Three bounded attempts, in order: Powell's hybrid method from the supplied
start; Levenberg--Marquardt from the same start if the first does not reach
the gate; and, when the caller passed a warm start, one more hybrid attempt
from the mode's own cold guess, since a warm start carried across a threshold
can land outside the basin. The best attempt is returned with its scaled
residual and its status, whether or not the gate was met.

\textbf{Non-convergence is a return value, never an exception.} The model is
used inside parameter scans that walk into regions where no positive-density
solution exists, and a scan must be able to score such a point and move on. A
malformed \emph{call} is a different thing --- an unknown mode, a fraction the
mode does not take, a sector the model does not implement --- and raises
before any solve, because it is a programming error that would otherwise be
repeated a million times in silence.

\subsection{Cold starts}
\label{sec:guesses}

Where the composition is unknown, the densities are estimated first and the
potentials follow from them. For the two $\beta$-equilibrium modes,
%
\begin{equation}
n_u = n_d = \nB , \qquad
n_s = \nB \cdot \min\!\Big(0.9,\; \max\big(0.01,\;
      \tfrac{T}{100\ \mathrm{MeV}} + \tfrac{\nB}{0.5\ \mathrm{fm}^{-3}}
      \big)\Big) ,
\label{eq:guess:beta}
\end{equation}
%
the strange fraction rising with both temperature and density because the
threshold of Eq.~(\ref{eq:kF}) is what suppresses it, and both loosen it. The
potentials are then $\mu_q = \sqrt{k_F^2 + m_q^2} + V$ at the Fermi momentum
of those densities,
%
\begin{equation}
k_F(n) = \hbar c \left( \frac{6 \pi^2 n}{g} \right)^{\!1/3} ,
\end{equation}
%
with $V$ from Eq.~(\ref{eq:vectorfield}) evaluated on the estimate, and
$\mu_e = \max(0, \mu_d - \mu_u)$, which is the $\beta$-equilibrium estimate
$-\muC$. The trapped mode adds $\munue = 10$~MeV.

Where the mode \emph{fixes} the composition, the densities are not estimated
but inverted from the constraints. At fixed $Y_C$ and $Y_S$,
%
\begin{equation}
n_s = Y_S \nB , \qquad
n_u = (1 + Y_C)\, \nB , \qquad
n_d = 3\nB - n_u - n_s ,
\label{eq:guess:ycys}
\end{equation}
%
which solves $R_4$, $R_5$ and $R_6$ of Eq.~(\ref{eq:res:ycys}) exactly and
leaves only the vector self-consistency for the solver to close. At fixed
$Y_C$ alone the strangeness is not determined by the constraints, and the
guess takes
%
\begin{equation}
n_s = 0.3\, \nB , \qquad
n_u = \max\big( \nB + Y_C \nB + \tfrac13 n_s,\; 0.3\, \nB \big) , \qquad
n_d = \max\big( \nB - \tfrac12 Y_C \nB,\; 0.3\, \nB \big) ,
\label{eq:guess:yc}
\end{equation}
%
the floors keeping a flavour from being seeded at or below zero. Where the
layout carries $\mu_e$ it is seeded at $\sqrt{k_{F,e}^2 + m_e^2}$ with
$k_{F,e} = \hbar c (3\pi^2 n_e)^{1/3}$ and $n_e = Y_C \nB$, which is the
neutralizing electron density.

A guess is only valid within its own mode, since the layouts of
Table~\ref{tab:unknowns} differ.

\subsection{Warm starts and the density sweep}

Along a density sweep each solved point seeds the next, in the layout of its
own mode: the potentials and the flavour densities of the previous solution,
read back into Table~\ref{tab:unknowns}'s order. They are the right
continuation variables because they vary smoothly with $\nB$ --- including
across the strange quark's onset, where the composition changes fastest --- and
because the large vector shift is common to all three potentials and therefore
cancels out of the step.

\section{The API surface}
\label{sec:api}

Three entry points, the same three every model in this repository exposes,
with the parameters always the first argument and the mode always required:
%
\begin{itemize}
\item \texttt{eos\_point(par, mode, species, n\_B=, T= \textbar\ SnB=,
      leptons=, **conditions)} --- one state. Returns a \texttt{PointResult}
      carrying \texttt{ok}, a \texttt{message}, and the \texttt{point} itself
      when \texttt{ok}. Exactly one of \texttt{T} and \texttt{SnB} must be
      given; \texttt{SnB} adds an outer solve for $T$. $\nB \le 0$ raises
      rather than returning a status: the equations have roots there --- a
      net-antiquark state solves them --- so the domain is stated rather than
      discovered.
\item \texttt{eos\_table(par, mode, species, axes, fixed=, leptons=,
      progress=, verbose=)} --- a solved grid over
      $\{\nB\} \times \{T\} \times$ the fraction axes the mode fixes, with the
      density axis warm-started. \texttt{skip\_errors=True}, the default,
      drops points the solver could not reach instead of aborting the table.
      \texttt{progress} is called once per completed line with
      $\{$\texttt{mode}, \texttt{line}, \texttt{n\_lines}, \texttt{temp\_key},
      \texttt{temp}, \texttt{fracs}, \texttt{n\_solved}, \texttt{n\_requested},
      \texttt{elapsed\_s}$\}$, the same dictionary in every model;
      \texttt{verbose=True} installs the built-in printer. Deep solver code
      never prints.
\item \texttt{eos\_response(par, mode, species, frozen='equilibrium', n\_B=,
      T=, leptons=, rel\_step=, **conditions)} --- second derivatives.
\end{itemize}

\subsection{What \texttt{eos\_response} implements}

Only \texttt{frozen='equilibrium'}: everything re-equilibrates under the
perturbation, so the derivatives are taken along the mode's own sequence. Both
are central differences over a relative step \texttt{rel\_step} (default
$10^{-3}$) in the variable differentiated, because vMIT has no analytic
Jacobian in this repository:
%
\begin{equation}
c_{s,\rm eq}^2 = \left.\frac{\partial P}{\partial \varepsilon}\right|_{T}
   \simeq \frac{P(\nB + \delta) - P(\nB - \delta)}
               {\varepsilon(\nB + \delta) - \varepsilon(\nB - \delta)} ,
\qquad
C_V = \frac{T}{\nB} \left.\frac{\partial s}{\partial T}\right|_{\nB}
   \simeq \frac{T}{\nB}\,
     \frac{s(T + \delta_T) - s(T - \delta_T)}{2\delta_T} ,
\label{eq:response}
\end{equation}
%
returned under the keys \texttt{cs2\_isothermal} and \texttt{C\_V}, the
latter only at $T > 0$ where it is defined. The adiabatic speed, larger by
$C_P/C_V$ at $T > 0$, is not computed by this model: $C_P$ is not among the
returned quantities, so there is no factor to form it with. Every other freeze of the specification --- frozen
per-species composition, frozen conserved fractions, the leptonic
re-neutralization variants, and the susceptibility matrix
$\chi_{ab} = \partial n_a / \partial \mu_b$ --- raises
\texttt{NotImplementedError} naming the gap (\S\ref{sec:deferred}).

\section{What a solved point returns}
\label{sec:returns}

A solve returns an \texttt{EoSPoint} carrying every quantity below. Nothing is
omitted because nothing downstream consumes it.

\begin{table}[h]
\centering
\begin{tabular}{lll}
\toprule
field & symbol & unit / meaning \\
\midrule
\texttt{converged} & --- & the status a caller must test first \\
\texttt{error}     & --- & largest scaled residual, Eq.~(\ref{eq:scales}),
                           dimensionless \\
\midrule
\texttt{n\_B}      & $\nB$        & $\mathrm{fm}^{-3}$, the condition \\
\texttt{T}         & $T$          & MeV, the condition \\
\texttt{Y\_C}, \texttt{Y\_S}, \texttt{Y\_L}
                   & $Y_C, Y_S, Y_{L_e}$ & the fractions the mode fixed \\
\midrule
\texttt{mu\_u}, \texttt{mu\_d}, \texttt{mu\_s}
                   & $\mu_u, \mu_d, \mu_s$ & MeV, \emph{physical} potentials
                                             (not $\mueff{q}$) \\
\texttt{mu\_e}, \texttt{mu\_nu}
                   & $\mu_e, \munue$ & MeV \\
\texttt{mu\_B}, \texttt{mu\_C}, \texttt{mu\_S}
                   & $\muB, \muC, \muS$ & MeV, Eq.~(\ref{eq:potentials}) \\
\midrule
\texttt{n\_u}, \texttt{n\_d}, \texttt{n\_s}
                   & $n_u, n_d, n_s$ & $\mathrm{fm}^{-3}$, net flavour
                                       densities \\
\texttt{n\_e}, \texttt{n\_nu}
                   & $n_e, n_{\nu_e}$ & $\mathrm{fm}^{-3}$ \\
\midrule
\texttt{P\_total}  & $P$           & $\mathrm{MeV\,fm^{-3}}$,
                                     Eq.~(\ref{eq:grandtotals}) \\
\texttt{e\_total}  & $\varepsilon$ & $\mathrm{MeV\,fm^{-3}}$ \\
\texttt{s\_total}  & $s$           & $\mathrm{fm}^{-3}$ \\
\midrule
\texttt{Y\_u}, \texttt{Y\_d}, \texttt{Y\_s}, \texttt{Y\_e}, \texttt{Y\_nu}
                   & $n_i/\nB$     & per-species fractions \\
\bottomrule
\end{tabular}
\caption{The fields of \texttt{EoSPoint}. When \texttt{converged} is
         \texttt{False} every other field holds the best iterate reached,
         which is not a physical state.}
\label{tab:returns}
\end{table}

Two of these are worth stating separately because \S11 of the repository
conventions singles them out.

\textbf{$s$} is \texttt{s\_total} of Eq.~(\ref{eq:grandtotals}), assembled as
the sum of the sectors' entropy densities. Each sector's own $s$ comes through
Eq.~(\ref{eq:sq}), the single-gas Euler identity at its own potential; the
model does not compute $s$ from the \emph{total} Euler relation, which is why
Eq.~(\ref{eq:euler}) is available as an independent check rather than being
true by construction.

\textbf{$n_s$} is the strange-quark number density, not a scalar density.
vMIT has no scalar field and no effective mass, so the trace identity
$n_s = (\varepsilon - 3P)/m^*$ is not the definition here and is not used
anywhere in the package; see \S\ref{sec:single}.

\subsection{A table row}

\texttt{eos\_table} flattens each solved point into a row keyed exactly the
way a hadronic table is keyed, so a pure-quark table and a hadronic one
concatenate without renaming: \texttt{n\_B}, \texttt{T}, \texttt{chi},
\texttt{phase}, \texttt{P}, \texttt{eps}, \texttt{s}, \texttt{S\_per\_B},
\texttt{mu\_B}, \texttt{mu\_C}, \texttt{mu\_S}, \texttt{mu\_e},
\texttt{Y\_C}, \texttt{Y\_S}, \texttt{Y\_u}, \texttt{Y\_d}, \texttt{Y\_s},
\texttt{Y\_e}, plus \texttt{Y\_nue} and \texttt{mu\_nue} in the trapped mode.
Here $\chi = 1$ and \texttt{phase}~$= $~\texttt{"Q"} say that the matter is
entirely deconfined. \texttt{Y\_C} in the row is the charge the solved state
turned out to have, computed from Eq.~(\ref{eq:charges}) --- not the value
that was requested, which is what makes the column meaningful in every mode.

\section{The \texttt{verify/} suite}
\label{sec:verify}

Eight physics invariants, each returning a structured pass/fail with the
largest error it saw:
%
\begin{enumerate}
\item \textbf{Euler relation}, Eq.~(\ref{eq:euler}), in every mode, at
      $10^{-8}$ relative.
\item \textbf{Free energy}, $f = \varepsilon - Ts = -P + \sum_i \mu_i n_i$, at
      $10^{-8}$.
\item \textbf{Vector self-consistency}: $V$ equals
      Eq.~(\ref{eq:vectorfield}) evaluated on the solved densities, and
      $n_q(\mu_q - V, T, m_q)$ reproduces the density the solver settled on,
      at $10^{-8}$.
\item \textbf{Bag and vector signs}: $\varepsilon_B = -P_B = B/\hc$ and
      $\varepsilon_V = +P_V$, at $10^{-14}$.
\item \textbf{Mode closures}: each mode's own conditions evaluated at its own
      solution --- $\muC + \mu_e$, $\muS$ and $\nC - n_e$ for
      \texttt{beta\_eq\_neutrinoless}; $\nC - Y_C\nB$ and $\muS$ for
      \texttt{fixed\_YC}; $\nC - Y_C\nB$ and $\nS - Y_S\nB$ for
      \texttt{fixed\_YC\_YS}; $\muC + \mu_e - \munue$ and the lepton fraction
      for the trapped mode --- at $10^{-8}$.
\item \textbf{Free-gas limit}: with $a = 0$ and $B = 0$ the solved state
      equals the Fermi integrals evaluated at the physical potentials
      themselves, at $10^{-10}$. This is the check that the interaction terms
      are additions to a correct kinetic sector rather than a correction
      hiding one.
\item \textbf{Residual gate}: every state the suite solved is inside
      $10^{-10}$.
\item \textbf{Causality}: $0 \le c_s^2 \le 1$ along a cold $\beta$-equilibrium
      sequence, with $c_s^2$ from Eq.~(\ref{eq:response}).
\end{enumerate}
%
The suite's density grid stays above the bag's unbinding point, where quark
matter exists at all: below it there is no positive-pressure solution to
check.

Not in this list, and deliberately: the monotonicity and causality gate a
table owes before it reaches a structure solver. vMIT's \texttt{table.py}
produces rows for a hybrid construction, where a pure quark branch is not the
delivered table; the gate belongs to whoever assembles the delivered one,
which for a hybrid star is \texttt{eos/mixed}.

\section{What is not implemented}
\label{sec:deferred}

Each of these raises, naming itself; none is silently ignored. The ledger is
\texttt{docs/DEFERRED.md}.

\begin{itemize}
\item \textbf{The muon lepton family.} \texttt{SpeciesFlags(muons=True)}
      raises, and \texttt{beta\_eq\_neutrino\_trapped} takes
      $(\nB, Y_{L_e}, T)$ only --- a $Y_{L_\mu}$ condition is refused at the
      API boundary.
\item \textbf{\texttt{thermal\_neutrinos}}, the flavours not tracked in the
      composition carried as $\mu = 0$ gases, is not wired and raises.
\item \textbf{The hadronic sector flags} \texttt{hyperons}, \texttt{deltas}
      and \texttt{thermal\_mesons} raise, and will continue to: they have no
      meaning in a deconfined phase, where strangeness enters through the $s$
      quark. This is a refusal for physics, not a gap.
\item \textbf{The entropy axis of a table.} \texttt{eos\_point} takes
      \texttt{SnB}; \texttt{TableSpec} raises for
      \texttt{axes=\{'SnB': ...\}}. The outer solve exists in
      \texttt{eos.general.tabulate}; wiring it into the table driver is what
      is left.
\item \textbf{Every response freeze but \texttt{equilibrium}}, and the
      susceptibility matrix $\chi_{ab}$. An analytic Jacobian is
      straightforward for this model --- one algebraic mean field, no scalar
      sector --- and would give $\chi_{ab}$ with it.
\item \textbf{Positivity of the flavour densities is not imposed.} At exotic
      fixed fractions ($Y_C$ well above $1$, say) the equations have solutions
      with net \emph{anti}-down and anti-strange densities, and the solver
      returns them as converged. They are genuine states of the model at
      finite temperature rather than solver failures, but nothing in the API
      says so.
\end{itemize}

\begin{acknowledgments}
Written as the specification of \texttt{eos/vmit}.
\end{acknowledgments}

\bibliography{../../docs/eos}

\end{document}
